\section{Metodología}

La práctica se desarrolló en GNU Radio Companion a partir del flujograma base suministrado en la guía del laboratorio \cite{guia_lab3,gnuradio}. El procedimiento se dividió en cuatro etapas.

\subsection{Señal binaria aleatoria bipolar}
Se generó una secuencia binaria aleatoria y se convirtió de unipolar a bipolar para reducir la componente DC. Luego se aplicó un bloque \textit{Interpolation FIR Filter} y se evaluaron los casos \(Sps=1\), \(4\), \(8\) y \(16\), registrando la forma temporal, la PSD y los parámetros \(R_b\), \(F_s\) y ancho de banda.

\subsection{Caracterización de ruido blanco}
Se configuraron las fuentes virtuales en los puntos \(p4\) y \(p5\) para observar el ruido blanco en el dominio del tiempo y en frecuencia. Además, se comparó el efecto de variar la amplitud del generador sobre el nivel de potencia espectral.

\subsection{Análisis con fuentes reales}
Se reemplazó el bloque \texttt{Random Source} por una cadena formada por \texttt{File Source}, \texttt{Unpack K Bits} y \texttt{Char to Float}, usando como entradas los archivos \texttt{rana.jpg} y \texttt{sonido.wav}. Con ello se comparó la PSD de fuentes reales frente a la de una fuente aleatoria ideal.

\subsection{Preguntas de control}
Finalmente, se modificó el vector de coeficientes \(h\) del filtro interpolador para evaluar cambios en codificación de línea y conformación de pulso, respondiendo de forma sintética las preguntas de control del laboratorio.